\chapter{Shards and Logical Processors}

CharlotteOS separates two concepts that are often conflated: the hardware
execution context (logical processor, or LP) and the userspace ownership domain
(shard). This chapter explains both and the relationship between them.

\section{The logical processor (LP) --- the hardware unit of scheduling}

An LP is CharlotteOS's representation of a schedulable hardware execution
context. In most configurations, one LP equals one CPU core (or one hyperthread).
The kernel shards its own state per-LP, meaning each LP has:

\begin{itemize}
    \item Its own run queue of ready threads.
    \item Its own timer queue and quantum timer.
    \item Its own scheduler instance (e.g., \texttt{RoundRobin}).
    \item Its own IPI command queues for receiving work from other LPs.
    \item Per-LP data structures accessible through \rusttype{PerLp<T>}.
\end{itemize}

An LP is a \textbf{place}: threads are pinned to LP at creation and (by default)
stay there. Migration is opt-in and restricted.

\section{The shard --- the userspace unit of ownership}

A shard is a userspace ownership and execution domain. The rules of the shard are
the sitas discipline:

\begin{enumerate}
    \item \textbf{Only the owning shard mutates shard-local state.} There is no
        lock on shard-local data because there is no concurrent access. The
        executor runs one task at a time on the shard.

    \item \textbf{Cross-shard interaction goes through bounded typed messages.}
        A value that moves from shard A to shard B is sent through a typed
        channel (\rusttype{ShardSender<M>}). The sender loses the value; the
        receiver gains it.

    \item \textbf{Placement is explicit.} Code that needs to be shard-local is
        accessed through \rusttype{ShardLocal<T>}, which verifies at runtime that
        the calling thread belongs to the owning LP.

    \item \textbf{Observability returns owned snapshots.} Diagnostics and
        introspection capture state by value. No borrowed reference into shard
        internals escapes the shard.
\end{enumerate}

\section{The mapping: one shard $\leftrightarrow$ one LP}

In the typical configuration:

\begin{verbatim}
one sitas shard
    <-> one LP-affine userspace thread
    <-> normally one logical CPU
\end{verbatim}

The userspace thread is pinned to its LP. Its executor runs only on that LP. Its
shard-local state (\rusttype{ShardLocal<T>}) is valid only when the owning LP's
thread is executing.

The architecture retains the distinction (shard $\neq$ LP) so that more complex
configurations remain possible:

\begin{itemize}
    \item \textbf{Multiple sharded processes:} one LP might host shards from
        multiple protection domains, time-sliced by the kernel scheduler.
    \item \textbf{CPU hotplug:} an LP can be taken offline; its shards can be
        migrated (when migration-safe) or torn down.
    \item \textbf{Oversubscription:} a test can run $N$ shards on $M < N$ LPs
        to exercise migration and scheduling paths.
\end{itemize}

\section{\texttt{PerLp<T>} --- interrupt-safe shared data}

\label{sec:per-lp}

For kernel data that must be accessible from multiple LPs (including from
interrupt context), CharlotteOS uses \rusttype{PerLp<T>}: a boxed slice of
per-LP cells, each wrapped in a lock. Access is:

\begin{verbatim}
per_lp_data.with(|lp_data| { lp_data.do_something(); });
\end{verbatim}

The \texttt{with} closure acquires the lock for the current LP, calls the closure,
and releases the lock. This is interrupt-safe (the lock masks interrupts during
the critical section on architectures where that is necessary).

Use \rusttype{PerLp<T>} for data that is:
\begin{itemize}
    \item Accessed from interrupt handlers (timers, device IRQs).
    \item Modified by an IPI from another LP.
    \item Part of the scheduler, memory allocator, or other kernel subsystem
        that must be LP-safe.
\end{itemize}

\section{\texttt{ShardLocal<T>} --- lock-free single-owner data}

\label{sec:shard-local}

For kernel data that is accessed only from one LP and never from interrupt
context, CharlotteOS provides \rusttype{ShardLocal<T>}: a lock-free per-LP
container using \texttt{UnsafeCell<T>} with runtime owner-check and borrow-flag
guard. Access is:

\begin{verbatim}
shard_local.with(|val| { *val = 42; });
\end{verbatim}

The \texttt{with} closure:
\begin{enumerate}
    \item Checks that the calling thread belongs to the owning LP (panics if not).
    \item Checks that the borrow flag is clear (panics on re-entrant access).
    \item Sets the borrow flag.
    \item Calls the closure with \texttt{\&mut T}.
    \item Clears the borrow flag.
\end{enumerate}

The closure must not:
\begin{itemize}
    \item Store the \texttt{\&mut T} reference anywhere (it must not escape).
    \item Send it to another thread or LP.
    \item Cross an \texttt{.await} point.
    \item Panic or unwind (the borrow flag would remain set, deadlocking future
        access).
\end{itemize}

\section{Why the distinction matters for critical software}

\begin{enumerate}
    \item \textbf{Data-race freedom by construction.} Shard-local state is never
        accessed concurrently. There is no lock, no atomic, no memory-ordering
        subtlety. If you hold \texttt{\&mut T}, no one else can access \texttt{T}.
        This is Rust's borrow-checker guarantee extended to the kernel's
        scheduling model.

    \item \textbf{Placement is deliberate.} Sending work to a
        specific LP is an explicit operation (\texttt{try\_run\_on\_lp},
        \texttt{ShardSender::try\_send}) and it is not at the mercy of a global
        work-stealing scheduler. This makes it possible to reason about locality,
        cache behavior, and latency.

    \item \textbf{The two container types force a design choice.} When you reach
        for shared data, you must decide: is this accessed from interrupt context?
        Use \rusttype{PerLp<T>} with locking. Is it thread-only? Use
        \rusttype{ShardLocal<T>} with zero overhead. There is no
        ``just wrap it in Mutex and move on'' ambiguity.

    \item \textbf{Cross-shard communication is typed and bounded.} A
        \rusttype{ShardSender<M>} carries a specific Rust type (not
        \texttt{Box<dyn Any>}, not a byte buffer). The receiver knows what it will
        get. The channel has fixed capacity; senders handle backpressure. There
        is no unbounded message queue that can exhaust memory.
\end{enumerate}

\endinput
